\chapter*{}
%\thispagestyle{empty}
%\cleardoublepage

%\thispagestyle{empty}

\input{portada/portada_2}



\cleardoublepage
\thispagestyle{empty}

\begin{center}
{\large\bfseries Sistema de gestión de logística de mercancías: aplicación BlockTrack}\\
\end{center}
\begin{center}
Juan José Martínez Águila\\
\end{center}

%\vspace{0.7cm}
\noindent{\textbf{Palabras clave}: Blockchain, trazabilidad, transporte de mercancías, logística, contrato inteligente,
aplicación web, transparencia, ethereum, sistema de seguimiento,...}\\

\vspace{0.7cm}
\noindent{\textbf{Resumen}}\\

En un mundo cada vez más interconectado, la gestión del transporte de mercancías se enfrenta al reto de 
garantizar la confianza, la transparencia y la eficiencia en cada etapa del proceso. 
A diario, miles de productos recorren largas distancias, pasando por manos de distintos agentes y 
atravesando múltiples puntos de control. Sin embargo, en muchos casos, 
la información sobre su recorrido sigue siendo opaca, fragmentada o difícil de verificar.\

De esta necesidad nace BlockTrack, una plataforma concebida para ofrecer una trazabilidad clara y confiable 
de cada envío. Su propósito es acompañar a las empresas y transportistas en la gestión diaria de sus operaciones, 
permitiendo registrar cada movimiento de manera segura y verificable gracias a la tecnología blockchain.\

Más que una simple herramienta digital, BlockTrack busca ser un espacio de conexión entre quienes intervienen
 en la cadena logística: enviadores, transportistas y destinatarios. A través de una interfaz web sencilla 
 y moderna, cada usuario puede seguir el estado de sus remesas en tiempo real, conocer su historial y 
 tener la certeza de que los datos no pueden ser alterados.\

El proyecto combina innovación tecnológica y sentido práctico, integrando la transparencia de blockchain 
con la accesibilidad de una aplicación web pensada para el uso cotidiano. En definitiva, 
BlockTrack aspira a transformar la forma en que entendemos la trazabilidad, 
convirtiéndola no solo en un proceso técnico, sino en una experiencia de confianza compartida.\
\cleardoublepage


\thispagestyle{empty}


\begin{center}
{\large\bfseries Goods Logistics Management System: BlockTrack Application}\
\end{center}
\begin{center}
Juan José Martínez Águila\\
\end{center}

%\vspace{0.7cm}
\noindent{\textbf{Keywords}: Blockchain, traceability, freight transport, logistics, smart contract, web application, transparency, ethereum, tracking system,...}\

\vspace{0.7cm}
\noindent{\textbf{Abstract}}\\

In an increasingly interconnected world, the management of freight transport faces the challenge of ensuring trust, transparency, and efficiency at every stage of the process. Every day, thousands of products travel long distances, passing through the hands of different agents and crossing multiple control points. However, in many cases, the information about their journey remains opaque, fragmented, or difficult to verify.\

From this need arises BlockTrack, a platform designed to provide clear and reliable traceability for every shipment. Its purpose is to support companies and carriers in the daily management of their operations, allowing each movement to be recorded securely and verifiably thanks to blockchain technology.\

More than just a digital tool, BlockTrack seeks to be a space of connection between those involved in the logistics chain: senders, carriers, and recipients. Through a simple and modern web interface, each user can track the status of their shipments in real time, review their history, and be certain that the data cannot be altered.\

The project combines technological innovation with practical purpose, integrating the transparency of blockchain with the accessibility of a web application designed for everyday use. Ultimately, BlockTrack aims to transform the way we understand traceability, turning it not only into a technical process but into an experience of shared trust.\

\chapter*{}
\thispagestyle{empty}

\noindent\rule[-1ex]{\textwidth}{2pt}\\[4.5ex]

Yo, \textbf{Juan José Martínez Águila}, alumno de la titulación \textbf{Ingeniería Informática} de la \textbf{Escuela Técnica Superior
de Ingenierías Informática y de Telecomunicación de la Universidad de Granada}, con DNI 20532681Y, autorizo la
ubicación de la siguiente copia de mi Trabajo Fin de Grado en la biblioteca del centro para que pueda ser
consultada por las personas que lo deseen.

\vspace{6cm}

\noindent Fdo: Juan José Martínez Águila

\vspace{2cm}

\begin{flushright}
Granada a 11 de Noviembre de 2025.
\end{flushright}


\chapter*{}
\thispagestyle{empty}

\noindent\rule[-1ex]{\textwidth}{2pt}\\[4.5ex]

D. \textbf{Nombre Apellido1 Apellido2 (tutor1)}, Profesor del Área de XXXX del Departamento YYYY de la Universidad de Granada.

\vspace{0.5cm}


\textbf{Informa:}

\vspace{0.5cm}

Que el presente trabajo, titulado \textit{\textbf{Sistema de gestión de logística de mercancías: aplicación BlockTrack}},
ha sido realizado bajo su supervisión por \textbf{Juan José Martínez Águila }, y autorizo la defensa de dicho trabajo ante el tribunal
que corresponda.

\vspace{0.5cm}

Y para que conste, expiden y firman el presente informe en Granada a 11 de Noviembre de 2025.

\vspace{1cm}

\textbf{El director:}

\vspace{5cm}

\noindent \textbf{Nombre Apellido1 Apellido2 (tutor1)}

\chapter*{Agradecimientos}
\thispagestyle{empty}

       \vspace{1cm}


En primer lugar, quiero agradecer a mis padres y a Ana por acompañarme en este proceso. Sin su apoyo incondicional y motivación, no habría sido posible completar este proyecto. 
Gracias por creer en mí y por enseñarme el valor del esfuerzo y la perseverancia.\\

También quiero expresar mi gratitud a los amigos que he conocido en la universidad, quienes han hecho que esta etapa sea inolvidable. Su compañía y apoyo han sido fundamentales para superar los desafíos académicos y personales.\\

Por último, agradecer a todos mis seres queridos que, de una u otra forma, han depositado su confianza en mí y
me han impulsado a seguir adelante. Este logro es tan suyo como mío. ¡Gracias a todos!
